\documentclass{article}
\usepackage[utf8]{inputenc}
\usepackage{amsmath}
\usepackage{graphicx}
\usepackage{geometry}
\usepackage{float}
\usepackage{booktabs}
\geometry{a4paper, margin=1in}

\title{Homework 6 - Part 2: Compressibility Effects in Slider Bearings}
\author{}
\date{}

\begin{document}
\maketitle

\section*{Part 2a: Compressible Reynolds Equation}
In this section, we account for the compressibility of air as peak pressures rise significantly above ambient conditions. Under the isothermal assumption, the fluid density is proportional to the absolute pressure $p_a = p + p_{atm}$. The resulting compressible Reynolds equation is nonlinear:

\begin{equation}
    -\nabla \cdot \left( \frac{p_a h^3}{12\mu} \nabla p \right) + \frac{1}{2} \nabla \cdot (uhp) = -\frac{p_{atm}}{2} \nabla \cdot (uh)
\end{equation}

where $p$ is the relative pressure and $p=0$ on the boundary $\partial \Omega$.

\subsection*{Iterative Solution Methodology}
The nonlinearity in the stiffness matrix $\bar{A}(p_a)$ was treated using a successive iteration strategy. Starting with an initial guess of $p_a = p_{atm}$, the system was iterated using:
\begin{equation}
    R(\bar{A}(p_a) - \bar{C})R^T \underline{p} = R\bar{C}\underline{\bar{p}}_{atm}
\end{equation}
The solution reached convergence (maximum pressure change $< 10^{-3}$) in 13 iterations.

\subsection*{Results}
The following parameters and computed values characterize the compressible taper-flat configuration:

\begin{itemize}
    \item \textbf{Trailing height ($h_2$):} $2.500 \times 10^{-7}$ m
    \item \textbf{Pitch ($\gamma$):} $2.927 \times 10^{-5}$ radians
    \item \textbf{Computed Load ($F$):} $0.27828$ N
    \item \textbf{Center of Pressure ($x_p$):} $2.194 \times 10^{-3}$ m
\end{itemize}

\section*{Pressure Distribution}
The absolute pressure distribution for the compressible case is shown in Figure \ref{fig:part_2a}. 

\begin{figure}[H]
    \centering
    \includegraphics[width=0.8\textwidth]{part_2a}
    \caption{Pressure distribution for the compressible taper-flat slider configuration.}
    \label{fig:part_2a}
\end{figure}

\end{document}